\section{Regras gerais}

\subsection{Mestre de Jogo}
O Mestre de Jogo (MJ) orienta o jogo e tem a última palavra. Mesmo que também interprete uma personagem, pode esclarecer ou ajustar uma regra quando necessário. O MJ resolve dúvidas, acompanha a pontuação quando se aplica e preenche as lacunas que o material escrito não cobre.

\subsection{Personagem}
Cada jogador interpreta uma personagem durante o jogo. A ficha dá-lhe informações e incentivos, em vez de um guião que tenha de seguir palavra por palavra.

As ações fora da personagem são proibidas, exceto quando um componente do jogo as exigir. Por exemplo, se quiseres revelar um dos teus segredos, deves fazê-lo através da interpretação, sem mostrar a tua ficha.

\ObjectivesHeading{Objetivos}
Os objetivos são resultados que dão uma recompensa à personagem que os alcançar. Alguns podem ser revelados apenas durante o jogo. A recompensa por cumprir um objetivo costuma ser indicada por \PositiveReward{1} e, se houver uma penalização por falhar, será indicada por \NegativeReward{1}.

\LosingConditionsHeading{Condições de derrota}
As condições de derrota são resultados que excluem uma personagem da classificação final, mesmo que tenha cumprido objetivos que valem pontos.

\SecretsHeading{Segredos}
Os segredos são factos que os outros jogadores podem descobrir.

Se um segredo apresentar a pontuação na forma (\NegativeReward{X}, \PositiveReward{Y}), o primeiro valor é a penalização que a personagem recebe se pelo menos outro jogador descobrir o segredo (\NegativeReward{X}); o segundo é a recompensa para os jogadores que o descobrirem (\PositiveReward{Y}).

Se um segredo tiver a forma (\PositiveReward{X}, \NegativeReward{Y}), é falso. Por cada outro jogador que o indicar como segredo desta personagem, ela recebe \PositiveReward{X}. Cada um desses jogadores recebe a penalização \NegativeReward{Y}.

Mesmo que não exista um segredo com a forma (\PositiveReward{X}, \NegativeReward{Y}), se um jogador atribuir a esta personagem um segredo que contradiga totalmente todos os seus segredos verdadeiros, esse jogador recebe \NegativeReward{1} e o jogador desta personagem recebe \PositiveReward{1}.

Os segredos só são avaliados no fim do jogo, quando os jogadores podem comunicar ao MJ \textbf{5} segredos que acreditam ter descoberto.

\ActionsHeading{Ações}
As ações são comportamentos ou tarefas que o jogador deve realizar, ou pelo menos ponderar seriamente, durante o jogo. Dão vida à interpretação e ajudam a fazer avançar a história.

\BackgroundHeading{Contexto}
O contexto apresenta o passado da personagem antes do evento. Pode incluir enquadramento narrativo, pistas para o mistério principal, informações que outros jogadores procuram e factos que a personagem quer esconder.

\subsection{Como jogar}
O grupo tenta resolver o mistério principal, que deverá tornar-se suficientemente claro para ser respondido antes de terminar a sessão.

Normalmente, isto culmina numa votação final, em que os jogadores têm de expulsar uma personagem vilã.

Mantém-te na personagem tanto quanto possível, evitando ações fora da interpretação. Trata o contexto como real, segue as tuas ações e leva os teus segredos a sério. Não leias a ficha em voz alta nem reveles informações muito específicas do jogo, exceto para pedir esclarecimentos ao MJ. Usa uma mensagem privada ou uma conversa discreta quando o MJ precisar de resolver algo em segredo.

Os pontos incentivam-te a cumprir objetivos secundários, a controlar as informações que partilhas e a investigar as outras personagens. Não percas de vista as condições de derrota. Não faz sentido proteger os teus segredos ao ponto de perder o jogo. Quando um segredo é exposto, a personagem afetada tem menos motivos para o esconder e pode tornar-se mais proativa.

Não te esqueças de que mentir é uma opção. Os jogadores podem revelar informações úteis enquanto se protegem com uma explicação alternativa plausível.

Também pode haver pistas escondidas nos espaços físicos, por isso os jogadores devem procurar e interagir com o cenário.

No fim do jogo, depois da votação, cada jogador terá a oportunidade de dizer ao MJ, em privado, \textbf{5} segredos que considera verdadeiros. Só depois se contam os pontos.
